% !TEX program = xelatex
% CUMCM 论文正式写作框架（电子版）
% 使用说明：复制为项目 06-paper/main.tex 后只修改项目副本；提交前清除占位符并逐页检查最终 PDF。
% 结构、内容和无状态交付执行 paper-writing.md，年度格式执行 OFFICIAL-CUMCM-001，所有结论必须来自可核验证据。

\documentclass[UTF8,a4paper,12pt,zihao=-4]{ctexart}

% ---------- 页面、字体与数学字体 ----------
\usepackage[top=2.5cm,bottom=2.5cm,left=2.5cm,right=2.5cm]{geometry}
\usepackage{fontspec}
\usepackage{unicode-math}

% Windows 优先使用宋体和 Times New Roman；其它系统使用相近的可移植后备字体。
% 发布前必须检查最终 PDF 字体属性，确认中文为宋体、英文与数字为 Times 系列。
\IfFontExistsTF{Times New Roman}
  {\setmainfont{Times New Roman}}
  {\setmainfont{TeX Gyre Termes}}
\IfFontExistsTF{SimSun}{
  \setCJKmainfont[AutoFakeBold=2.5]{SimSun}
  \setCJKsansfont[AutoFakeBold=2.5]{SimSun}
  \setCJKmonofont{SimSun}
}{
  \setCJKmainfont[AutoFakeBold=2.5]{FandolSong-Regular}
  \setCJKsansfont[AutoFakeBold=2.5]{FandolSong-Regular}
  \setCJKmonofont{FandolFang-Regular}
}
\IfFontExistsTF{Consolas}{\setmonofont{Consolas}}{\setmonofont{Latin Modern Mono}}
\setmathfont{TeX Gyre Termes Math}

% ---------- 公式、图表、列表、代码与引用 ----------
\usepackage{amsmath}
\usepackage{graphicx}
\usepackage{subcaption}
\usepackage{booktabs,longtable,tabularx,array,multirow}
\usepackage{enumitem}
\usepackage{listings}
\usepackage{xcolor}
\usepackage{fancyhdr}
\usepackage{hyperref}
\usepackage[nameinlink,noabbrev]{cleveref}

% ---------- 全文样式 ----------
\setlength{\parindent}{2em}
\setlength{\parskip}{0pt}
% 小四正文按约 20 pt 基线距排印，对应写作规范的紧凑 1.5 倍视觉节奏。
\linespread{1.38}
\setlist{nosep,leftmargin=2em}
\setlength{\emergencystretch}{2em}

\ctexset{
  section={
    name={,、},
    number=\chinese{section},
    format=\centering\bfseries\zihao{-3},
    beforeskip=1.5ex,
    afterskip=1.0ex
  },
  subsection={
    format=\raggedright\bfseries\zihao{-4},
    beforeskip=1.2ex,
    afterskip=0.8ex
  },
  subsubsection={
    format=\raggedright\bfseries\zihao{-4},
    beforeskip=1.0ex,
    afterskip=0.6ex
  }
}

\pagestyle{fancy}
\fancyhf{}
\fancyfoot[C]{\thepage}
\renewcommand{\headrulewidth}{0pt}
\renewcommand{\footrulewidth}{0pt}

\captionsetup{
  font=small,
  labelfont=bf,
  labelsep=quad,
  justification=centering,
  singlelinecheck=true
}
\captionsetup[subfigure]{font=small,labelfont=normalfont,labelsep=space}

\hypersetup{
  hidelinks,
  unicode=true,
  pdfauthor={},
  pdfsubject={},
  pdfkeywords={},
  pdftitle={CUMCM 论文}
}

\crefname{figure}{图}{图}
\crefname{table}{表}{表}
\crefname{equation}{式}{式}
\crefname{section}{节}{节}

\lstset{
  basicstyle=\ttfamily\small,
  numbers=left,
  numberstyle=\ttfamily\scriptsize,
  numbersep=8pt,
  frame=single,
  breaklines=true,
  breakatwhitespace=false,
  columns=fullflexible,
  keepspaces=true,
  showstringspaces=false,
  tabsize=4,
  captionpos=t
}

% ---------- 模板变量与安全占位 ----------
\newcommand{\TemplateField}[1]{\textbf{【#1】}}
\newcommand{\PaperTitle}{\TemplateField{填写准确、具体的论文题目}}

% 根据实际情况设置 AI 使用状态。使用本模板生成或修改论文时不得作虚假声明。
\newif\ifusedai
\usedaitrue
\newcommand{\AIUsageSummary}{语言润色、代码调试与排版核验等}

% 根据实际情况设置程序和支撑材料状态；三者必须与附录及提交压缩包一致。
\newif\ifusedprogram
\usedprogramtrue
\newif\ifhassupportmaterials
\hassupportmaterialstrue

% 图片缺失时仍可编译，但会显示醒目的占位框；正式提交前不得保留占位框。
\newcommand{\safeincludegraphics}[2][]{%
  \IfFileExists{#2}{%
    \includegraphics[#1]{#2}%
  }{%
    \fbox{%
      \begin{minipage}[c][4cm][c]{0.82\linewidth}
        \centering
        \textbf{图像占位}\par
        \texttt{\detokenize{#2}}\par
        正式提交前必须替换为实际图片
      \end{minipage}%
    }%
  }%
}

% 从支撑材料源文件直接排印完整代码；缺失时显示阻断提示。
\newcommand{\includecode}[3]{%
  \IfFileExists{#3}{%
    \lstinputlisting[language=#1,caption={#2}]{#3}%
  }{%
    \par\noindent
    \fbox{%
      \begin{minipage}{0.94\linewidth}
        \textbf{代码缺失：}\texttt{\detokenize{#3}}。正式提交前必须补齐完整、可运行且与论文结果一致的源程序。
      \end{minipage}%
    }
    \par
  }%
}

\begin{document}
\pagenumbering{arabic}
\setcounter{page}{1}

% ============================================================
% 摘要内容与版式执行 paper-writing.md 和 OFFICIAL-CUMCM-001。
% ============================================================
\thispagestyle{fancy}
\begin{center}
  {\bfseries\zihao{3}\PaperTitle\par}
  \vspace{1.2em}
  {\bfseries\zihao{4}摘要\par}
\end{center}

\TemplateField{用两至三句话概括研究对象、总体建模主线和主要成果。}

针对问题一，\TemplateField{说明任务和关键转化}，建立\textbf{【核心模型一】}，采用\TemplateField{求解方法}得到\textbf{【关键数值、单位或明确结论】}；通过\TemplateField{检验方法及指标}验证结果可信。

针对问题二，\TemplateField{说明与前问的继承关系及新增任务}，构建\textbf{【核心模型二】}，得到\textbf{【最终方案、排序、预测或分类结果】}；\TemplateField{给出误差、稳健性或约束检查结论}。

针对问题三，\TemplateField{说明情景变化和模型调整}，使用\textbf{【核心模型三】}求得\textbf{【关键结果】}；\TemplateField{说明比较或敏感性分析结果}。

针对问题四，\TemplateField{说明开放任务或决策目标}，提出\textbf{【最终策略或方案】}；\TemplateField{说明适用范围、可信度和主要限制}。

综上，\textbf{【概括模型体系的主要价值、边界与可推广结论】}。

\vspace{0.8em}
\noindent\textbf{关键词：}\TemplateField{研究对象}；\TemplateField{核心模型}；\TemplateField{求解算法}；\TemplateField{验证方法}；\TemplateField{决策主题}

\clearpage

% ============================================================
% 正文、目录和页数执行 paper-writing.md 与 OFFICIAL-CUMCM-001。
% ============================================================
\section{问题重述}

\subsection{问题背景}
\TemplateField{用自己的语言概括背景、研究对象和应用意义，不照抄赛题，不提前展开模型和结果。}
相关背景事实、公开数据和专业判断必须就近引用经核验来源，例如\cite{topic-ref-1,topic-ref-2}。

\subsection{问题要求}
\begin{enumerate}[label=\arabic*.]
  \item \TemplateField{问题一：写清输入、条件、约束和输出。}
  \item \TemplateField{问题二：写清与问题一的依赖关系及新增输出。}
  \item \TemplateField{问题三：写清情景变化、约束和输出。}
  \item \TemplateField{问题四：写清最终决策、附件或评价要求。}
\end{enumerate}

\section{问题分析}

\subsection{问题一的分析}
\TemplateField{说明数学本质、数据或机理特征、核心困难、候选模型、选择理由、预期输出和验证路径。}

\subsection{问题二的分析}
\TemplateField{说明继承的变量、参数和结果，以及新增目标、约束和验证路径。}

\subsection{问题三的分析}
\TemplateField{说明情景变化如何映射到参数、指标或模型结构。}

\subsection{问题四的分析}
\TemplateField{说明开放任务如何转化为可计算、可评价的数学问题。}

% 保留的单图写法：图必须先在正文中引用，后出现，并在图后解释其作用。
总体建模流程如\cref{fig:overall-flow}所示，该图用于说明各子问题的数据和模型继承关系。
\begin{figure}[htbp]
  \centering
  \safeincludegraphics[width=0.72\textwidth]{figures/overall-flow.pdf}
  \caption{总体建模流程}
  \label{fig:overall-flow}
\end{figure}
\TemplateField{解释流程图中的关键路径，不得用流程图替代模型说明。}

\section{模型假设}
\begin{enumerate}[label=\textbf{假设\arabic*.}]
  \item \TemplateField{必要假设及其物理、业务或统计依据。}
  \item \TemplateField{必要假设及其合理性；不得用假设直接消除题目核心难点。}
  \item \TemplateField{对结果影响显著的假设，并在后文灵敏度或模型评价中回扣。}
\end{enumerate}

\section{符号说明}
% 符号表优先采用“符号—含义—单位”三列；局部临时变量可在公式附近定义。
\begin{table}[htbp]
  \centering
  \caption{主要符号及其含义}
  \label{tab:symbols}
  \begin{tabularx}{0.88\textwidth}{>{\centering\arraybackslash}p{0.18\textwidth}X>{\centering\arraybackslash}p{0.18\textwidth}}
    \toprule[1.2pt]
    符号 & 含义 & 单位 \\
    \midrule[0.8pt]
    $x_i$ & \TemplateField{第 $i$ 个核心变量的含义} & \TemplateField{单位或—} \\
    $p_j$ & \TemplateField{第 $j$ 个模型参数的含义} & \TemplateField{单位或—} \\
    $f(\symbf{x})$ & \TemplateField{目标函数或评价指标的含义} & \TemplateField{单位或—} \\
    \bottomrule[1.2pt]
  \end{tabularx}
\end{table}

\section{模型的建立与求解}

\subsection{问题一模型的建立与求解}

\subsubsection{问题一的建模准备}
\TemplateField{给出本问正式建模前必需的坐标、几何判定、数据变换、指标口径、预备公式、继承关系或输入输出；简单问题可用三至八行，复杂内容可配公式和示意图。}
模型抽象和算法设计可参考数学建模基础教材\cite{modelbook,algobook}，但具体关系必须由本题数据、机理和约束推导。

\subsubsection{问题一模型的建立}
\noindent\TemplateField{局部模型名称——说明本模型在问题一中的作用}
\TemplateField{依次定义变量和指标，推导局部关系，然后集中汇总完整模型。公式前写依据，公式后解释含义。}

\noindent\textbf{完整模型汇总：}\TemplateField{集中列出本问的核心方程、变量范围、参数条件、输入与输出；优化模型还须列出目标函数和全部约束。}

% 保留并规范的优化模型汇总写法：非优化问题不得机械使用目标函数和 s.t.。
若本问属于优化问题，可将完整模型集中写为
\begin{equation}
\label{eq:q1-model}
\begin{aligned}
\min_{\symbf{x}}\quad & f(\symbf{x}) \\
\mathrm{s.t.}\quad &
\left\{
\begin{aligned}
g_1(\symbf{x}) &\leq 0,\\
g_2(\symbf{x}) &\leq 0,\\
h_1(\symbf{x}) &= 0,\\
\symbf{x} &\in \Omega .
\end{aligned}
\right.
\end{aligned}
\end{equation}
\cref{eq:q1-model}中，$f(\symbf{x})$表示\TemplateField{目标的实际含义}，$g_1$、$g_2$和$h_1$分别描述\TemplateField{资源、业务或机理约束}，$\Omega$为\TemplateField{变量定义域}。正式论文应替换整组公式，不得保留无关示例。

\subsubsection{问题一模型的求解}
\TemplateField{写明求解方法、关键步骤、影响结果的参数或容差、约束与边界检查；复杂算法可配编号步骤。}

\subsubsection{问题一结果的分析}
先直接回答本问：\textbf{【填写问题一的最终数值、单位、方案或结论】}。

% 保留的三线结果表写法：表题在上，单位放入表头，小数位和有效数字统一。
主要结果见\cref{tab:q1-results}。
\begin{table}[htbp]
  \centering
  \caption{问题一主要结果}
  \label{tab:q1-results}
  \begin{tabular}{cccc}
    \toprule[1.2pt]
    对象 & 核心指标（单位） & 误差或状态（单位） & 最终结论 \\
    \midrule[0.8pt]
    \TemplateField{对象1} & \TemplateField{数值} & \TemplateField{数值或状态} & \TemplateField{结论} \\
    \TemplateField{对象2} & \TemplateField{数值} & \TemplateField{数值或状态} & \TemplateField{结论} \\
    \bottomrule[1.2pt]
  \end{tabular}
\end{table}
\TemplateField{解释表中主要发现、实际意义、判据和不确定性，不逐行朗读。}
\TemplateField{至少给出两类适用的结果分析、误差分析、模型检验、敏感性或等效验证。}

\subsection{问题二模型的建立与求解}

\subsubsection{问题二的建模准备}
\TemplateField{说明从问题一继承的变量、参数、结果或预备方法，并给出本问新增的坐标、变换、指标、约束、参数和输出。}

\subsubsection{问题二模型的建立}
\noindent\TemplateField{局部模型名称——说明本模型在问题二中的作用}
\TemplateField{完成“问题转化—变量指标—局部关系—完整模型汇总”，并引用经核验的题目相关方法来源\cite{topic-ref-3}。}

\noindent\textbf{完整模型汇总：}\TemplateField{集中列出问题二的目标函数或核心方程、约束或初边值条件、定义域、参数条件、输入与输出。}

% 保留的并排子图写法：各子图与总图均有唯一标签。
模型结构或方案比较如\cref{fig:q2-comparison}所示。
\begin{figure}[htbp]
  \centering
  \begin{subfigure}[t]{0.31\textwidth}
    \centering
    \safeincludegraphics[width=0.95\linewidth]{figures/q2-case-a.pdf}
    \caption{情景一}
    \label{fig:q2-case-a}
  \end{subfigure}\hfill
  \begin{subfigure}[t]{0.31\textwidth}
    \centering
    \safeincludegraphics[width=0.95\linewidth]{figures/q2-case-b.pdf}
    \caption{情景二}
    \label{fig:q2-case-b}
  \end{subfigure}\hfill
  \begin{subfigure}[t]{0.31\textwidth}
    \centering
    \safeincludegraphics[width=0.95\linewidth]{figures/q2-case-c.pdf}
    \caption{情景三}
    \label{fig:q2-case-c}
  \end{subfigure}
  \caption{问题二不同情景的结果比较}
  \label{fig:q2-comparison}
\end{figure}
\TemplateField{分别解释各子图差异，并说明它们如何支撑问题二结论。}

\subsubsection{问题二模型的求解}
\TemplateField{写明实际算法、参数选择、数据划分或求解器设置，以及约束、泄漏、收敛或边界检查。}

\subsubsection{问题二结果的分析}
先直接回答本问：\textbf{【填写问题二最终结果】}。\TemplateField{给出至少两类适用分析或等效验证。}

\subsection{问题三模型的建立与求解}

\subsubsection{问题三的建模准备}
\TemplateField{说明情景变化对应的变量、参数、边界或数据口径变化，以及从前问继承的预备关系。}

\subsubsection{问题三模型的建立}
\noindent\TemplateField{局部模型名称——说明本模型在问题三中的作用}
\TemplateField{定义新增变量与参数，推导并汇总完整模型；相关专业依据引用经核验来源\cite{topic-ref-4}。}

\noindent\textbf{完整模型汇总：}\TemplateField{集中列出问题三的目标函数或核心方程、约束或初边值条件、定义域、参数条件、输入与输出。}

\subsubsection{问题三模型的求解}
\TemplateField{写明求解流程、关键参数和可行性检查。}

\subsubsection{问题三结果的分析}
先直接回答本问：\textbf{【填写问题三最终结果】}。\TemplateField{报告敏感性、稳健性、误差或其它等效证据。}

\subsection{问题四模型的建立与求解}

\subsubsection{问题四的建模准备}
\TemplateField{给出开放任务正式建模前所需的评价口径、候选方案表示、输入输出、情景边界或前问继承关系。}

\subsubsection{问题四模型的建立}
\noindent\TemplateField{局部模型名称——说明本模型在问题四中的作用}
\TemplateField{给出变量、指标、局部关系与完整模型汇总，说明输出如何对应题目要求。}

\noindent\textbf{完整模型汇总：}\TemplateField{集中列出问题四的目标函数或核心方程、约束或初边值条件、定义域、参数条件、输入与输出。}

\subsubsection{问题四模型的求解}
\TemplateField{写明实际求解或方案生成过程。}

\subsubsection{问题四结果的分析}
先直接回答本问：\textbf{【填写问题四最终方案或结论】}。\TemplateField{说明验证、适用条件和主要限制。}

\section{模型的评价与推广}

\subsection{模型的优点}
\TemplateField{每项优点对应具体模型结构、运行结果或检验证据，不写空泛评价。}

\subsection{模型的缺点}
\TemplateField{如实说明假设、数据、参数、数值方法和适用边界的限制。}

\subsection{模型的改进与推广}
\TemplateField{说明推广对象、需调整的变量和约束、需补充的数据及推广后的验证方式。}

% ============================================================
% 2026 AI 新规：AI 工具使用声明必须放在参考文献之前，二选一。
% AI 声明及配套文件执行 OFFICIAL-CUMCM-001。
% ============================================================
\clearpage
\section*{AI 工具使用声明}
\ifusedai
本参赛队在竞赛过程中使用了 AI 工具，主要用于\AIUsageSummary，详细使用情况见支撑材料。
\else
本参赛队在竞赛过程中未使用任何 AI 工具。
\fi

% 文献组成、核验和正文引用执行 PW-CITE-001 与 WG-EVID-001。
\renewcommand{\refname}{七、参考文献}
\begin{thebibliography}{99}
  \bibitem{modelbook}
  姜启源, 谢金星, 叶俊.
  \newblock 数学模型（第五版）[M].
  \newblock 北京: 高等教育出版社, 2018.

  \bibitem{algobook}
  司守奎.
  \newblock 数学建模算法与应用[M].
  \newblock 北京: 国防工业出版社, 2011.

  \bibitem{topic-ref-1}
  \TemplateField{作者. 与题目背景直接相关的文献题名[类型]. 出版信息或期刊信息, 年份.}

  \bibitem{topic-ref-2}
  \TemplateField{作者或机构. 经核验的数据、标准或权威资料题名[类型]. 年份. URL/DOI及访问日期.}

  \bibitem{topic-ref-3}
  \TemplateField{作者. 与核心模型或算法直接相关的文献题名[类型]. 期刊或出版社, 年份. DOI.}

  \bibitem{topic-ref-4}
  \TemplateField{作者. 与参数、验证或应用场景直接相关的文献题名[类型]. 期刊或出版社, 年份. DOI.}
\end{thebibliography}

% ============================================================
% 附录：总体另起一页，每个独立附录再次另起一页。
% 支撑材料清单、实际压缩包、论文引用和代码必须完全一致。
% ============================================================
\clearpage
\section*{附录A\quad 支撑材料文件清单}
\ifhassupportmaterials
\begin{longtable}{>{\centering\arraybackslash}p{0.08\textwidth}p{0.34\textwidth}p{0.46\textwidth}}
  \caption{支撑材料文件清单}\label{tab:support-files}\\
  \toprule[1.2pt]
  序号 & 相对路径或文件名 & 内容及其与正文的对应关系 \\
  \midrule[0.8pt]
  \endfirsthead
  \multicolumn{3}{c}{续表\quad 支撑材料文件清单}\\
  \toprule[1.2pt]
  序号 & 相对路径或文件名 & 内容及其与正文的对应关系 \\
  \midrule[0.8pt]
  \endhead
  \midrule
  \multicolumn{3}{r}{续下页}\\
  \endfoot
  \bottomrule[1.2pt]
  \endlastfoot
  1 & \texttt{code/q1.py} & \TemplateField{问题一求解程序，对应式、表或图编号} \\
  2 & \texttt{code/q2.py} & \TemplateField{问题二求解程序，对应式、表或图编号} \\
  3 & \texttt{data/processed-data.csv} & \TemplateField{自主查阅或处理后的必要数据；不得重复提交赛题原始数据} \\
  \ifusedai
  4 & \TemplateField{按 OFFICIAL-CUMCM-001 填写} & \TemplateField{AI 使用配套文件及用途} \\
  \fi
  \TemplateField{序号} & \TemplateField{其余实际文件名} & \TemplateField{用途及正文对应关系} \\
\end{longtable}
\else
本论文没有支撑材料。
\fi

\clearpage
\section*{附录B\quad 建模所用完整源程序}
\ifusedprogram
% 正式提交前按实际文件逐一调用；不得只附伪代码、核心片段或无法运行的节选。
% 路径按“模板已复制到项目 06-paper/”设置，必要时按项目实际结构调整。
\subsection*{B.1\quad 问题一源程序}
本程序用于\TemplateField{说明所解决的子问题以及与正文公式、图表和结果的对应关系}。
\includecode{Python}{问题一完整源程序}{../03-models/q1.py}

\subsection*{B.2\quad 问题二源程序}
本程序用于\TemplateField{说明所解决的子问题以及与正文公式、图表和结果的对应关系}。
\includecode{Python}{问题二完整源程序}{../03-models/q2.py}

% 若存在问题三、问题四或绘图、数据处理程序，应继续逐一完整列出。
\else
本论文没有用到程序。
\fi

\end{document}
